% Image credits for the photographs and AI illustrations of Book 3 (University, Year 1), Indonesian edition.
% Machine-readable license records live in images/book3/CREDITS.md; keep this
% file in sync with frontmatter/image-credits-book3.tex (one \item per item there).
% Every credit, author, dataset and licence is kept exactly as in the English file.
\chapter*{Kredit Gambar}

Gambar, peta dan grafik dalam buku ini bersifat orisinal; peta dan grafik yang
dibuat dari data yang diterbitkan menyebutkan sumbernya dalam keterangan gambar
atau dalam daftar ini. Foto digunakan di bawah lisensi bebasnya masing-masing,
dengan terima kasih kepada para penulisnya; tautan sumber lengkap dan URL lisensi
untuk setiap foto tercantum dalam \texttt{images/book3/CREDITS.md} di repositori buku.

\bigskip
Selain foto dan gambar orisinal, beberapa figur adalah ilustrasi yang dihasilkan AI,
dibuat dengan pembuatan gambar OpenAI dari prompt yang ditulis oleh editor buku, dan
ditinjau keakuratan ilmiahnya sebelum dimasukkan. Gambar-gambar ini hanya menampilkan
adegan umum, tidak pernah tempat nyata, orang nyata, atau spesimen nyata. Prompt
lengkap untuk setiap gambar yang dihasilkan tercatat dalam
\texttt{images/book3/ai/PROMPTS.md} di repositori.

\bigskip
\noindent Foto dan data, menurut bab:
\begin{itemize}\small
  \item Bab~1 (astronomi sferis): jejak bintang di atas Observatorium Paranal, R.~Wesson/ESO, CC BY 4.0.
  \item Bab~2 (teleskop, detektor dan fotometri): enam segmen cermin primer teleskop luar angkasa James Webb sebelum pengujian kriogenik, NASA/MSFC/David Higginbotham, domain publik.
  \item Bab~3 (planetologi komparatif): planet-planet berbatu dalam skala, NASA; Venus dari radar Magellan, NASA/JPL; Mars dari wahana Viking 1, NASA/USGS; Europa, NASA/JPL/DLR; Io, NASA/JPL/University of Arizona; semuanya domain publik.
  \item Bab~4 (peta geologi): William Smith, \emph{A Delineation of the Strata of England and Wales, with part of Scotland}, 1815; domain publik (Wikimedia Commons).
  \item Bab~5 (penampang): ketidakselarasan Hutton di Siccar Point, PawełS, CC BY-SA 4.0 (Wikimedia Commons).
  \item Bab~6 (metode lapangan dan laboratorium): granit biotit dalam sayatan tipis di antara polar bersilang, Darla Sondrol, CC0 (Wikimedia Commons).
  \item Bab~7 (mineral pembentuk batuan): spesimen olivin, kuarsa, biotit, almandin dan hornblende, Robert M.~Lavinsky, CC BY-SA 3.0; labradorit, Twyla Baker, CC0; romboeder belahan kalsit, Giovanna Canu dan Eva Santini, CC BY 4.0 (semuanya Wikimedia Commons).
  \item Bab~8 (mineralogi optik): gabro dalam sayatan tipis dengan satu polar dan di antara polar bersilang, Chmee2, CC BY-SA 3.0; Auguste Michel-Lévy sebagai mahasiswa, domain publik (Wikimedia Commons). Warna interferensi dihitung dari fungsi pencocokan warna CIE 1931 dan iluminan standar CIE D65, \copyright{} CIE, CC BY-SA 4.0.
  \item Bab~9 (Bumi berlapis): Andrija Mohorovičić, domain publik; Inge Lehmann pada 1932, foto oleh Even Neuhaus, domain publik (Wikimedia Commons). PREM: A.~M.~Dziewonski dan D.~L.~Anderson (1981), seperti ditabulasikan dalam ObsPy.
  \item Bab~10 (anggaran panas Bumi): aliran panas global dari J.~H.~Davies dan D.~R.~Davies (2010), \emph{Solid Earth}, CC BY 3.0; adiabat mantel dari T.~Katsura (2022), \emph{JGR Solid Earth}, CC BY.
  \item Bab~11 (isostasi dan anomali gravitasi): alat bandul Vening Meinesz, Broodt, CC BY-SA 4.0 (Wikimedia Commons). Ketinggian dari ETOPO1 (NOAA National Centers for Environmental Information), diambil sampel setiap satu derajat.
  \item Bab~12 (kinematika lempeng pada bola): sebuah stasiun penentuan posisi satelit permanen, Alpsdake, CC BY-SA 4.0 (Wikimedia Commons). Batas lempeng dari PB2002 (P.~Bird, 2003; GeoJSON oleh H.~Ahlenius, ODC-By 1.0); garis pantai dari Natural Earth (domain publik).
  \item Bab~13 (rekahan, punggungan dan tepi pasif): retakan terbuka di Þingvellir, Islandia, Olga Ernst, CC BY-SA 4.0; lava bantal di Rekahan Galápagos, NOAA Okeanos Explorer Program, CC BY 2.0 (keduanya Wikimedia Commons). Kedalaman dari ETOPO1 (NOAA National Centers for Environmental Information) dan umur dasar samudra dari M.~Seton dan lainnya (2020), EarthByte, CC BY 4.0.
  \item Bab~14 (tepi aktif): Puncak Sarychev sedang meletus, dilihat dari Stasiun Luar Angkasa Internasional, NASA, domain publik (Wikimedia Commons). Gempa dari USGS ANSS Comprehensive Catalog (domain publik); gunung api dari Program Vulkanisme Global Smithsonian Institution; batas lempeng dari PB2002 (P.~Bird, 2003; GeoJSON oleh H.~Ahlenius, ODC-By 1.0); pantai dari Natural Earth (domain publik); umur dasar samudra dari M.~Seton dan lainnya (2020), EarthByte, CC BY 4.0.
  \item Bab~15 (titik panas dan vulkanisme intralempeng): air mancur lava di Pu\textquoteleft u \textquoteleft O\textquoteleft o, Kilauea, G.~E.~Ulrich, U.S. Geological Survey, domain publik; Ghats Barat di Matheran, Nicholas (Nichalp), CC BY-SA 2.5 (keduanya Wikimedia Commons). Posisi titik panas, azimut jejak dan laju dari Tabel~1 suplemen elektronik W.~J.~Morgan dan J.~P.~Morgan (2007); kedalaman dari ETOPO1 (NOAA National Centers for Environmental Information); batas lempeng dari PB2002 (P.~Bird, 2003; GeoJSON oleh H.~Ahlenius, ODC-By 1.0); pantai dari Natural Earth (domain publik).
  \item Bab~16 (tegangan, regangan dan reologi): sel triaksial tekanan tinggi untuk mengukur kekuatan batuan, Darryl Peroni, CSIRO, CC BY 3.0 (Wikimedia Commons). Parameter rayapan dislokasi kuarsit basah (Gleason dan Tullis, 1995) dan olivin kering (Hirth dan Kohlstedt, 2003) seperti dihimpun oleh M.~A.~Lopez-Sanchez dalam \emph{strength\_envelopes} (MPL-2.0).
  \item Bab~17 (sesar, lipatan dan fabrik): cermin sesar di Corona Heights, San Francisco, Dextersinister, CC BY-SA 4.0; breksi sesar Keystone Thrust, Nevada, Jstuby, domain publik; lipatan chevron di Millook, Cornwall, Smalljim, CC BY 3.0; antiklin dengan belahan bidang aksial dekat Potes, Cantabria, Mikenorton, CC BY-SA 3.0 (semuanya Wikimedia Commons).
  \item Bab~18 (sumber gempa): tanah dan pagar yang terpisah oleh rekahan 1906 di Marin County, \emph{Popular Science Monthly} vol.~69 (1906), domain publik; pergeseran Wallace Creek melintasi Sesar San Andreas, R.~E.~Wallace, U.S. Geological Survey, domain publik (keduanya Wikimedia Commons). Gempa dari USGS ANSS Comprehensive Catalog (domain publik); mekanisme fokus gempa Tohoku 2011 dari momen tensor fase-W USGS (domain publik).
  \item Bab~19 (pelapukan dan lempung): granit lapuk dan pasir granit di Targassonne, Pyrenees, Pinpin, CC~BY~2.5; bauksit pada batupasir kaolinitik di Weipa, Australia, Werner Schellmann, CC~BY-SA~2.5; kristal kaolinit di bawah mikroskop elektron pemindai, U.S. Geological Survey, domain publik (semuanya Wikimedia Commons). Jari-jari ion dari R.~D.~Shannon (1976).
  \item Bab~20 (erosi dan transpor): gumpalan banjir Sungai Mississippi dari instrumen MODIS pada Terra, 4~Maret 2018, citra NASA Earth Observatory oleh Joshua Stevens, domain publik (Wikimedia Commons).
  \item Bab~21 (sedimen klastik dan lingkungannya): perlapisan silang dalam Batupasir Aztec, Nevada, Chip Carroon, Bureau of Land Management, CC~BY-SA~2.0; delta Sungai Lena dari Landsat~7, U.S. Geological Survey, domain publik; sebuah turbidit dengan pembagian urutan Bouma, Becke-Oese, Jo Weber, CC~BY-SA~3.0 (semuanya Wikimedia Commons).
  \item Bab~22 (sedimen karbonat, silika dan evaporit): Bora Bora dari Stasiun Luar Angkasa Internasional, NASA Johnson Space Center, domain publik; sel \emph{Emiliania huxleyi}, Alison R.~Taylor, University of North Carolina Wilmington Microscopy Facility, CC~BY~2.5; \emph{Orbulina universa}, H.~Takagi, K.~Kimoto, T.~Fujiki, H.~Saito, C.~Schmidt, M.~Kucera dan K.~Moriya, CC~BY~4.0; diatom dari McMurdo Sound, Gordon T.~Taylor, Stony Brook University, NOAA Corps Collection, domain publik (semuanya Wikimedia Commons).
  \item Bab~23 (diagenesis dan cekungan sedimen): tumbuhan kuarsa dalam Batupasir Potsdam, M.~C.~Rygel, CC~BY-SA~4.0 (Wikimedia Commons). Tumpukan muka laut dari R.~M.~Spratt dan L.~E.~Lisiecki (2016), NOAA NCEI World Data Service for Paleoclimatology; parameter pemadatan dari J.~G.~Sclater dan P.~A.~F.~Christie (1980) seperti ditabulasikan dalam pyBacktrack (EarthByte).
  \item Bab~24 (skala stratigrafi): paku emas GSSP Selandian di Zumaia, Josugoni, CC~BY-SA~4.0 (Wikimedia Commons); batas Kapur--Paleogen di Trinidad Lake State Park, Nationalparks, CC~BY-SA~2.5 (Wikimedia Commons). Umur dasar dan warna Bagan Kronostratigrafi Internasional, International Commission on Stratigraphy, CC~BY~4.0.
  \item Bab~25 (metode penanggalan radiometrik): butir zirkon dalam katodoluminesensi, Emmanuel Roquette, CC~BY-SA~4.0; spektrometer massa ionisasi termal, Radiogenic, CC0 (keduanya Wikimedia Commons). Kurva kalibrasi IntCal20 dari P.~J.~Reimer dan lainnya (2020), \emph{Radiocarbon} 62, CC~BY~4.0, seperti didistribusikan dengan paket rintcal (M.~Blaauw).
  \item Bab~26 (struktur atmosfer dan samudra): pengambilan CTD dan rosette di Selat Bering, NASA Goddard Space Flight Center, domain publik (Wikimedia Commons). Suhu dan salinitas World Ocean Atlas 2023 (NOAA National Centers for Environmental Information).
  \item Bab~27 (radiasi dan transpor panas): Bumi yang terkena matahari dari Deep Space Climate Observatory, NASA Goddard Space Flight Center, domain publik (Wikimedia Commons). Radiasi, fluks permukaan dan presipitasi dari NCEP-NCAR Reanalysis~1 (E.~Kalnay dan lainnya, 1996), data disediakan oleh NOAA PSL, Boulder, Colorado, USA, dari situs webnya di https://psl.noaa.gov.
  \item Bab~28 (sirkulasi umum atmosfer): siklon ekstratropis di Teluk Alaska, NASA MODIS/LANCE, domain publik; zona konvergensi intertropis dari data awan GOES-11, NASA Earth Observatory, domain publik (keduanya Wikimedia Commons). Angin dan tekanan permukaan laut dari NCEP-NCAR Reanalysis~1 (E.~Kalnay dan lainnya, 1996), data disediakan oleh NOAA PSL, Boulder, Colorado, USA, dari situs webnya di https://psl.noaa.gov.
  \item Bab~29 (sirkulasi samudra): Arus Teluk yang terlihat dari suhu permukaan laut, NOAA NESDIS, domain publik; pelampung profil samudra dalam yang sedang diturunkan, fotografer NOAA yang tidak diketahui, domain publik (keduanya Wikimedia Commons).
\end{itemize}
